%!TEX root = ./main.tex

\section[Waveforms]{Where You Put the Symbols}

% SECTION TIMING: 22:00--33:00 (4 frames).
% ACCURACY NOTE: OTFS maps delay-Doppler symbols onto a time-frequency grid via the
% ISFFT, so it inherits an OFDM-like modulator. ODDM instead builds pulses that are
% orthogonal with respect to the delay-Doppler resolution directly. Do not describe
% ODDM as "OTFS with better pulses" -- the orthogonality claim is the whole point.

\begin{frame}{OFDM already carries a radar inside it}
  \vspace{0.3em}
  \begin{columns}[T,onlytextwidth]
    \column{0.44\textwidth}
      \centering
      \begin{tikzpicture}[scale=0.80,transform shape]
        \foreach \x in {0,...,5}
          \foreach \y in {0,...,3}
            \draw[draw=softgray!50,fill=blue!6] (\x*0.62+0.12,\y*0.52+0.10)
              rectangle ++(0.58,0.48);
        \draw[-{Latex[length=2mm]},softgray] (0,0) -- (4.0,0)
          node[anchor=north,font=\scriptsize] {time (symbols)};
        \draw[-{Latex[length=2mm]},softgray] (0,0) -- (0,2.4);
        \node[font=\scriptsize,rotate=90,anchor=south] at (-0.16,1.2) {frequency};
      \end{tikzpicture}
    \column{0.52\textwidth}
      \small
      \begin{itemize}\itemsep5pt
        \item Mature: Wi-Fi, 4G, 5G, 6G.
        \item The grid is already knobs.
        \item A cell serves bits or echoes.
        \item No new modulator needed.
      \end{itemize}
  \end{columns}

  \vspace{0.4em}
  \qa{If OFDM is good enough, why design new waveforms?}
     {Because OFDM is described in time and frequency, and a moving target is
      naturally described in delay and Doppler. Fast targets smear an OFDM grid;
      they sit still in a delay--Doppler one.}

  \glossline{OFDM = Orthogonal Frequency Division Multiplexing}
\end{frame}

\begin{frame}{Delay--Doppler is a different place to stand}
  \vspace{0.4em}
  \centering
  \begin{tikzpicture}[scale=0.88,transform shape]
    % time-frequency
    \draw[draw=ranblue,fill=blue!5] (0,0) rectangle (2.6,2.0);
    \node[font=\scriptsize,anchor=north] at (1.3,-0.05) {time};
    \node[font=\scriptsize,rotate=90,anchor=south] at (-0.12,1.0) {frequency};
    \node[font=\small\bfseries,text=ranblue,anchor=south] at (1.3,2.10) {Time--frequency};
    \foreach \p in {(0.5,1.5),(1.1,0.6),(2.0,1.2),(1.6,1.8),(0.8,0.3)}
      \fill[ranblue!70] \p circle (1.6pt);
    % arrow
    \draw[flow,draw=softgray,thick] (3.1,1.0) -- (4.7,1.0);
    \node[font=\scriptsize,anchor=south] at (3.9,1.10) {ISFFT};
    % delay-doppler
    \draw[draw=coreorange,fill=orange!6] (5.2,0) rectangle (7.8,2.0);
    \node[font=\scriptsize,anchor=north] at (6.5,-0.05) {delay};
    \node[font=\scriptsize,rotate=90,anchor=south] at (5.08,1.0) {Doppler};
    \node[font=\small\bfseries,text=coreorange,anchor=south] at (6.5,2.10) {Delay--Doppler};
    \fill[coreorange] (6.2,1.2) circle (2.4pt);
    \fill[coreorange] (7.0,0.8) circle (2.4pt);
  \end{tikzpicture}

  \vspace{0.9em}
  \Large
  A target that \textcolor{ranblue}{smears} across time and frequency\\[0.25em]
  is a \textcolor{coreorange}{point} in delay and Doppler.

  \glossline{ISFFT = Inverse Symplectic Finite Fourier Transform}
\end{frame}

\begin{frame}{Three waveforms, three bets}
  \vspace{0.2em}
  \scriptsize
  \setlength{\tabcolsep}{4pt}
  \renewcommand{\arraystretch}{1.30}
  \begin{tabularx}{\textwidth}{>{\bfseries}p{1.1cm}
      >{\raggedright\arraybackslash}p{2.1cm}
      >{\raggedright\arraybackslash}X
      >{\raggedright\arraybackslash}X}
    \toprule
    & \textbf{Symbols live in} & \textbf{Strength} & \textbf{Weakness} \\
    \midrule
    OFDM & Time--frequency & Deployed everywhere; subcarriers double as radar bins &
      Doppler breaks subcarrier orthogonality \\
    OTFS & Delay--Doppler, mapped up by ISFFT & Robust to high mobility; simple to
      bolt onto OFDM & Leakage between bins; high PAPR \\
    ODDM & Delay--Doppler, natively & Truly orthogonal pulses; sparse, clean echoes &
      Early research; hardware immature \\
    \bottomrule
  \end{tabularx}

  \glossline{OTFS = Orthogonal Time Frequency Space \quad
    ODDM = Orthogonal Delay-Doppler Division Multiplexing\\
    PAPR = Peak-to-Average Power Ratio, which decides how hard the amplifier works}

  \vspace{0.4em}
  \takeaway{OFDM is what ships.\\ODDM is what the sensing side wants.}
\end{frame}

\begin{frame}{Why orthogonality is worth the trouble}
  \vspace{0.5em}
  \begin{columns}[T,onlytextwidth]
    \column{0.47\textwidth}
      \centering
      {\small\textbf{\textcolor{softgray}{Leaky basis (OTFS)}}\par}
      \vspace{0.5em}
      \begin{tikzpicture}[scale=1.05,transform shape]
        \draw[draw=softgray!60] (0,0) rectangle (3.0,1.7);
        \fill[coreorange!85] (1.0,1.0) circle (2.8pt);
        \foreach \p in {(0.72,1.0),(1.28,1.0),(1.0,0.74),(1.0,1.26),(0.78,0.78),(1.22,1.22)}
          \fill[coreorange!35] \p circle (2.0pt);
        \fill[coreorange!85] (2.1,0.6) circle (2.8pt);
        \foreach \p in {(1.84,0.6),(2.36,0.6),(2.1,0.36),(2.1,0.84)}
          \fill[coreorange!35] \p circle (2.0pt);
      \end{tikzpicture}
    \column{0.47\textwidth}
      \centering
      {\small\textbf{\textcolor{softgreen}{Orthogonal basis (ODDM)}}\par}
      \vspace{0.5em}
      \begin{tikzpicture}[scale=1.05,transform shape]
        \draw[draw=softgray!60] (0,0) rectangle (3.0,1.7);
        \fill[coreorange!90] (1.0,1.0) circle (2.8pt);
        \fill[coreorange!90] (2.1,0.6) circle (2.8pt);
      \end{tikzpicture}
  \end{columns}

  \vspace{0.9em}
  \qa{Two blurred dots or two sharp dots --- why does the estimator care?}
     {Leakage is energy in the wrong bin. It raises the sidelobes, hides a weak
      target next to a strong one, and shows up as interference between users.}
\end{frame}
